
\section{Conclusion and Future Work}
\label{sec:conclusion}

This paper presents a mixed-methods research framework
for investigating textual markers of emotional attachment
and social isolation in accounts of interaction with
AI companions. The framework combines qualitative
content analysis, normalized marker frequencies,
non-parametric statistical tests, and regression analysis.

For RQ2, the synthetic pipeline test produced higher
emotional-dependence marker frequencies in the
long-duration group than in the short-duration group.
The statistical procedures successfully generated
group comparisons, effect-size estimates, and
confidence intervals. These results demonstrate
pipeline functionality but do not constitute
empirical evidence about human behavior.

For RQ3, the synthetic data yielded a weak positive
Spearman correlation between engagement duration
and social-avoidance indicators. However, the adjusted
OLS model did not produce a statistically significant
duration coefficient and explained less than 1\%
of the outcome variation. These findings illustrate
the importance of comparing statistical approaches
and interpreting their results cautiously.

Future work will focus on three priorities.
First, an authentic corpus will be collected and
analyzed with documented ethical and privacy safeguards.
Second, author-level aggregation and coding reliability
will be verified using the planned codebook and
independent coders. Third, the analysis will be
extended to additional platforms and populations,
with sensitivity analyses to evaluate the robustness
of the observed associations.

Overall, the current contribution is a reproducible
research design and a tested analytical workflow,
rather than confirmed psychological findings
about AI-companion users.
